% Source: https://artofproblemsolving.com/community/c3h2586378_advanced_linear_equations_prealgebra
% Stable identity: aops-topic-2586378-q06
If $a=c$, the equation reads $0\cdot x=d-b$.
\begin{itemize}
\item If $b\ne d$, the right side is nonzero, while the left side is $0$ for every $x$: there are no solutions.
\item If $b=d$, both sides are $0$ for every $x$: the two sides of the original equation are the same expression $ax+b$, and \emph{every} real number is a solution.
\end{itemize}
So $a=c$ alone only says the formula fails; it does not decide between ``no solutions'' and ``infinitely many''. That is why $b\ne d$ is also needed.
